% Component derivation for inclusion in the separate machine notes.
% This file is intentionally not a replacement for the manuscript.
\subsection{The Ising decomposition}

We keep the auxiliary Ramond parity on each cut.  The two irreducible
Virasoro modules in the Neveu--Schwarz Fock space have highest weights
$0$ and $1/2$.  Each Ramond parity gives an irreducible module of highest
weight $1/16$.  The field inserted on the additional sphere has weight
$1/2$.  The central charge in this subsection is that of the Ising
Virasoro algebra, namely $c=1/2$; it is not the superconformal central
charge.

The auxiliary three-point forms in the spin frames of the human notes
obey
\begin{align}
 \rho_{\mathsf F}(\mathbf1,u^0,u^0)&=1,&
 \rho_{\mathsf F}(\mathbf1,u^1,u^1)&=i,\\
 \rho_{\mathsf F}(\psi_{-1/2}\mathbf1,u^0,u^1)&=-\frac1{\sqrt2},&
 \rho_{\mathsf F}(\psi_{-1/2}\mathbf1,u^1,u^0)&=\frac{i}{\sqrt2}.
\end{align}
All other ground components vanish by parity.  To obtain the second
line, move the fermion contour at infinity onto the Ramond puncture at
one.  The Ward relation gives $i\rho_{\mathsf F}
(\psi_{-1/2}\mathbf1,u^{\mathsf b},u^{1-\mathsf b})+
\rho_{\mathsf F}(\mathbf1,\psi_0u^{\mathsf b},u^{1-\mathsf b})=0$.
Substitution of $\psi_0u^{\mathsf b}=u^{1-\mathsf b}/\sqrt2$ gives
the displayed values.

The ordered middle insertion acts on a ground state as
\begin{equation}
 Q\Theta\psi_0u^{\mathsf b}
 =\frac{Q}{\sqrt2}(-1)^{\mathsf b}u^{\mathsf b}.
\end{equation}
It commutes with parity and is a primary intertwiner of Ising weight
$1/2$, since $\Theta$ commutes with every auxiliary Virasoro generator.
Consequently its descendant matrix elements are precisely those in the
Virasoro one-point block.  Combining the two outer ground forms with
this middle matrix element and the theta sewing signs gives
the following sign ledger.  The middle raw pairing is $Q/\sqrt2$
in each row; the signs from its outgoing ground norm have already
been included.  The norm column is the product of the four internal
ground norms, and the last sign includes the NS fermion and theta
reordering signs:
\[
\begin{array}{c|cc|ccc}
 h_{\NS}&\mathsf b&\mathsf c&\rho_{\mathsf F}^{2}
 &\text{norm product}&\text{sewing sign}\\\hline
 0&0&0&1&1&1\\
 0&1&1&-1&-1&-1\\
 \tfrac12&0&1&\tfrac12&1&1\\
 \tfrac12&1&0&-\tfrac12&-1&1
\end{array}
\]
Thus the complete Ising decomposition is
\begin{align}
 \mathbb F_{\mathsf F}[\Theta v_{1/2}]
 ={}&\frac{Q}{\sqrt2}(1-\eta_2\eta_3)
 F\left(0,\frac1{16},\frac1{16},\frac1{16};\frac12;\frac12
 \mathrel\big|q_1,\hat q_{2,1},\hat q_{2,2},q_3\right)\nonumber\\
 &+\frac{Q}{2\sqrt2}(\eta_1\eta_2+\eta_1\eta_3)q_1^{1/2}
 F\left(\frac12,\frac1{16},\frac1{16},\frac1{16};\frac12;\frac12
 \mathrel\big|q_1,\hat q_{2,1},\hat q_{2,2},q_3\right).
 \label{eq:ising-ordered-theta-decomposition}
\end{align}
Both blocks on the right must use the \emph{irreducible} Ising modules.
Their primary coefficients are one.  The four internal entries refer,
in order, to the Neveu--Schwarz edge, the two parts of the split Ramond
edge, and the remaining Ramond edge.  The next entry is the external
weight and the last is the central charge.  The $q_1^{1/2}$ in the
second term restores the physical level of the Neveu--Schwarz fermion.
This formula contains the complete parity dependence; all remaining
coefficients are ordinary scalar Virasoro coefficients.
